% !TEX root = ../main.tex

\chapter{相关技术与理论基础}

\section{股票短期走势预测任务概述}

股票短期走势预测关注的是在给定历史信息和当前市场状态后，判断未来较短时间窗口内的价格方向或收益变化。相比中长期价值分析，短期走势更容易受到事件冲击、投资者情绪、市场关注度和信息扩散速度影响。只使用历史价格与成交量，通常难以解释这些短期扰动。新闻文本、社交媒体情绪和其他非结构化信息已被多项研究用于股票预测任务，大语言模型也开始被用于从金融文本中提取市场相关线索\cite{bollen2011twitter,zhang2025gnn,zhou2025expertopinion,lopezlira2023chatgptstock,elahi2024financialnewsllm}。

从任务形式上看，股票短期走势预测既可以建模为收益率回归问题，也可以建模为涨跌方向分类问题。考虑到真实市场中的短期收益分布噪声较大、极易受到偶然因素干扰，而方向判断更符合实际交易决策中的“买入—观望—卖出”逻辑，本文将任务定义为面向未来 1 日、3 日和 7 日的涨跌方向预测任务。该任务不仅需要模型理解历史价格所反映的市场状态，还需要综合评论、新闻、研报以及实时 Web 检索结果等外部信息，判断当前时点下市场叙事、风险事件和投资预期可能对未来走势产生的影响。

与传统时间序列模型相比，基于大语言模型的股票短期预测更依赖语义建模。模型不仅要处理价格序列中的统计规律，还要从多源异构文本中抽取事件、观点、情绪、因果线索和潜在共识，再将这些信息映射为未来价格方向判断。因此，本文任务并不是单纯的行情预测，而是同时涉及金融文本分析、检索增强生成、监督微调与强化学习对齐等多个方向\cite{lewis2020rag,ye2024r2ag,lee2024finllms,li2023llmfinance,nie2024llmfinancialapplications,yang2023fingpt,wu2023bloomberggpt}。

\section{多源金融信息的特点与作用}

金融文本来源差异很大。评论更新快但噪声高，新闻更接近事件传播，研报偏向专业分析，Web 检索结果适合补充临时热点和交叉验证。本文没有把这些材料简单合并成一段长输入，而是先按来源分别处理，再送入最终预测环节。这样做牺牲了一部分原文细节，但能减少不同质量、不同时间尺度的信息互相干扰\cite{lee2024finllms,li2023llmfinance,nie2024llmfinancialapplications}。

\subsection{评论信息的时效性与情绪性}

评论信息主要来自股吧、论坛、社交媒体和投资社区，其特点是更新快、情绪强、噪声大。在市场出现热点题材、突发消息或板块异动时，评论通常比正式新闻更早反映投资者关注点和态度变化。本文保留评论分支，是因为它可能捕捉到新闻尚未充分展开的短期情绪；但它不能直接等同于预测信号。Bollen 等研究发现，社交媒体中的群体情绪与市场变化之间存在统计关联，为评论型文本进入金融预测提供了早期证据\cite{bollen2011twitter}。

但评论信息并不等同于有效信号。其内容中往往夹杂大量重复表达、情绪宣泄、非理性观点甚至误导性信息。因此，如果直接将评论原文输入模型，容易放大噪声干扰。后续研究进一步指出，社交媒体和论坛的真正价值不在于原始帖子数量本身，而在于能否识别出持续提供有效判断的用户、代表性观点及其扩散关系。基于图结构的情绪分析方法和专家意见挖掘方法均说明，对评论信息进行筛选、归纳和结构化处理，比简单做情绪正负分类更能提升预测价值\cite{zhang2025gnn,zhou2025expertopinion}。

在本文框架中，评论主要用于补充价格序列之外的市场情绪和题材热度。若直接把原始评论输入模型，重复口号和情绪宣泄可能被放大；若完全丢弃评论，又会损失散户讨论热度这一类短期信号。因此，系统选择先归纳近期讨论主题、多空分歧和可能的情绪倾向，再将其作为辅助材料使用。

\subsection{新闻信息的事件驱动特征}

与评论信息相比，新闻信息具有更强的事件驱动属性。公告、政策、行业动态、公司经营变化、国际局势和突发风险事件，往往首先通过新闻媒体被市场广泛感知，并进一步影响投资者预期。短期股价波动中，很多显著变动都与特定事件直接相关，因此新闻文本是短期预测中不可忽视的重要来源。已有研究表明，大语言模型能够从新闻标题和报道中识别对未来收益具有解释力的关键信息，这说明新闻并不仅仅提供事实，还能承载影响价格方向的市场叙事与因果线索\cite{lopezlira2023chatgptstock,elahi2024financialnewsllm}。

新闻的可信度和结构化程度通常高于评论，事件边界也更清楚。不过，新闻可能滞后于最早的市场讨论，不同媒体的叙述角度也会不同，单条新闻的影响强度还依赖当时的行情背景。本文把新闻视为“事件触发型证据”，主要用于帮助模型识别短期走势中的外部驱动因素。

\subsection{研报信息的专业分析特征}

研报信息主要来自券商、研究机构或专业分析团队，其突出特点是专业术语多、逻辑链条完整、分析视角相对系统。与评论和新闻相比，研报通常更关注行业逻辑、公司基本面、估值判断、政策影响和未来催化剂，因此在解释“为什么涨”“为什么跌”方面具有更强的分析深度。金融领域相关研究和综述普遍指出，金融文本具有鲜明的领域术语和上下文依赖，普通通用模型往往难以充分理解其中的专业含义，这也是金融专用模型和金融领域微调持续受到重视的重要原因\cite{araci2019finbert,lee2024finllms,li2023llmfinance,nie2024llmfinancialapplications,wu2023bloomberggpt}。

研报信息的局限在于覆盖面和时效性并不总是最强，而且不同机构的观点可能带有自身研究框架和立场偏好。因此，研报更适合作为“高质量分析证据”来增强模型的逻辑判断能力，而不是作为唯一信息来源。对于本文任务而言，研报的主要价值在于为模型提供更专业的行业解释和未来预期分析，从而在评论和新闻之外增加一层较高质量的推理依据。

\subsection{Web 检索结果的补充作用}

在真实预测场景中，单纯依赖预先收集的静态数据往往不足以反映市场最新变化。尤其是在个股或板块出现临时热点、政策变化或舆情波动时，模型若无法接入实时外部信息，就很容易受到参数化知识陈旧性的限制。RAG 研究表明，通过在生成过程中接入外部检索文档，可以显著提升模型在知识密集任务中的事实性与更新能力\cite{lewis2020rag}。R\textsuperscript{2}AG 则进一步指出，仅仅“检索到文档”还不够，还需要通过更合理的提示设计和检索信息组织方式，帮助模型真正理解并利用检索结果\cite{ye2024r2ag}。

在本文框架中，Web 检索结果承担补充和校验作用。它能补足评论、新闻和研报未覆盖到的新事实，也能为单一来源信息提供交叉验证。遇到新事件、新概念或临时热点时，检索结果可以扩展模型可见的上下文范围，因此被纳入在线预测流程。

\section{大语言模型基本原理}

本文使用大语言模型，主要是因为股票分析材料大多不是干净的表格。评论里有情绪和口语化表达，新闻里有事件线索，研报里有较长的分析链条，搜索结果又会带来来源和时间上的差异。传统分类器可以处理单一字段，但要在同一个问题里比较这些文本，大语言模型更方便作为统一的阅读和归纳模块。

从模型结构看，现代大语言模型以 Transformer 为基础。Transformer 通过自注意力机制建模序列中不同位置之间的关系，使模型能够在较长上下文中捕捉前后依赖\cite{vaswani2017attention}。本文并不展开 Transformer 的完整推导，而是利用这一类模型已经具备的长文本理解和指令执行能力，把不同来源的股票材料组织成自然语言输入，让模型输出可检查的判断。

金融场景对模型提出了额外要求。BloombergGPT、FinGPT 以及相关综述都指出，金融文本术语密集、时效性强，且对事实约束比较敏感\cite{lee2024finllms,li2023llmfinance,nie2024llmfinancialapplications,yang2023fingpt,wu2023bloomberggpt}。因此，在本文系统中，基础模型不能只依赖预训练知识，还需要结合实时抓取材料、明确的 Prompt 约束以及后续微调。换言之，大语言模型在这里不是直接替代交易策略，而是承担材料阅读、观点归纳和辅助判断的角色。

\section{基于 Prompt 的金融文本理解方法}

在本文中，Prompt 首先是工程上的输入规范。评论、新闻、研报和网页结果的写法差异很大，如果直接拼接给模型，模型很容易只写一段泛泛总结，或者忽略某一类材料。Prompt 的作用就是告诉模型：每一类材料应该看什么，哪些信息需要保留，最后回答应当落到股票走势分析上。

金融文本对 Prompt 也比较敏感。同样一条新闻，既可以被总结成普通事件，也可以被分析为对公司经营、市场情绪或短期价格预期的影响。已有研究表明，大语言模型能够从新闻标题和金融报道中抽取与收益方向相关的信息，但前提是任务定义和输入组织足够清楚\cite{lopezlira2023chatgptstock,elahi2024financialnewsllm}。FinGPT 等工作也说明，金融任务往往需要任务指令和轻量微调配合，而不是只依赖通用模型的自由生成\cite{yang2023fingpt}。

因此，本文的 Prompt 不只是“提问模板”。分组阶段先把长文本压短，预测阶段再把不同来源放到同一个判断问题中。这样既能控制上下文长度，也便于检查模型是否依据新闻、研报或检索结果给出结论。该思路与 RAG 和 R\textsuperscript{2}AG 中“先组织外部证据，再引导模型使用证据”的做法一致\cite{lewis2020rag,ye2024r2ag}。

\section{监督微调方法}

本文需要监督微调，原因并不是基础模型完全不能回答股票问题，而是它的回答容易不稳定。未微调模型有时会解释很多背景，却不按要求给出最终标签；有时会把“上涨、下跌、震荡”等表达混在一起，给自动评测带来困难。SFT 的作用，是让模型熟悉本文的数据格式、回答顺序和标签约束。

监督微调通常使用“输入—输出”样本继续训练已有模型，相比从头训练成本更低，也更适合金融领域这种标注数据有限但任务格式明确的场景\cite{yang2023fingpt,ouyang2022instructgpt}。FinBERT 和 FinGPT 等研究表明，领域适配和任务微调可以改善金融文本理解效果\cite{araci2019finbert,yang2023fingpt}。本文选择 LoRA 而非全参数微调，主要考虑到 7B 模型的训练资源开销较大，同时希望保留基础模型原有的指令能力，只通过 adapter 快速适配股票方向预测格式。

在本文任务中，SFT 更像是把模型“调到题目格式上”。训练样本要求模型先给出简短分析，再输出可解析的“涨”或“跌”。这样做不会保证预测一定正确，但可以减少格式漂移，让后续离线评测和 GRPO 奖励计算有稳定入口。

\section{强化学习对齐方法（GRPO）}

SFT 解决的是“会不会按格式作答”的问题，但它不直接优化最终评价指标。对于本文的方向预测任务，模型可能已经能输出三项标签，却不一定在不同候选答案中选择更准确、更一致的那一个。因此，SFT 之后继续考虑强化学习，主要是为了把“预测是否命中、格式是否可解析、理由是否与结论一致”这些目标写进训练信号中。InstructGPT 也说明，在监督微调后引入偏好或奖励信号，可以让模型输出更贴近任务要求\cite{ouyang2022instructgpt}。

如果采用传统 PPO，通常还需要额外训练价值模型来估计基线，这会增加显存占用和调试复杂度\cite{schulman2017ppo}。本文的训练环境和任务规模都不适合再引入一套复杂 critic，因此选择 GRPO 作为后续方案。GRPO 来自 DeepSeekMath，其做法是在同一输入上采样多个候选输出，再用组内相对得分构造优化信号，从而避免显式训练独立 critic\cite{shao2024deepseekmath}。

这一点和本文任务比较贴合。同一段行情输入可以生成多种分析过程和方向组合，系统可以根据真实标签、输出格式和解析结果给这些候选答案打分。得分高的回答在组内获得更大权重，模型就有机会逐渐偏向更准确、更稳定的输出。本文引入 GRPO，并不是把股票预测说成完整交易策略，而是在 SFT 已经跑通后，为后续继续按评价目标优化模型留下可执行路径。

\section{本章小结}

本章的重点不是重新讲一遍通用模型原理，而是说明后续系统为什么需要这些工具。短期股价判断里，行情只能给出已经发生的价格状态，评论、新闻、研报和检索结果则补充事件、情绪和预期。大语言模型负责阅读和归纳这些材料，Prompt 负责限定输入与输出，SFT 让回答格式稳定，GRPO 则为后续按评价目标继续优化留下接口。
